% !TEX root = Vicente_Calisto_Resume_EN.tex
%-------------------------------------------------------------------------------
% Resume content, English and Spanish together.
%
% Every piece of prose is written as \tr{English}{Español}, so changing one
% language means editing the other right next to it. Build both PDFs with
% `make` (see README.md). Keep it to ONE page: CI fails the build otherwise.
%-------------------------------------------------------------------------------

\begin{document}

%----------HEADING----------
\begin{center}
    {\Huge \scshape Vicente Calisto} \\ \vspace{3pt}
    \small \tr{Senior Software Engineer $\cdot$ Database Engines $\cdot$ C++}
              {Ingeniero de Software Senior $\cdot$ Motores de Bases de Datos $\cdot$ C++}
    \\ \vspace{3pt}
    Chile $|$
    \href{mailto:vicente.calisto98@gmail.com}{\underline{vicente.calisto98@gmail.com}} $|$
    \href{https://linkedin.com/in/vicente-calisto}{\underline{linkedin.com/in/vicente-calisto}} $|$
    \href{https://github.com/VicenteVicente}{\underline{github.com/VicenteVicente}}
\end{center}

%----------EXPERIENCE----------
\section{\tr{Experience}{Experiencia}}
  \resumeSubHeadingListStart

    \resumeSubheading
      {\tr{Millennium Institute for Foundational Research on Data (IMFD)}{Instituto Milenio Fundamentos de los Datos (IMFD)}}
      {Santiago, Chile}
      {\tr{Software Engineer, Database Engines}{Ingeniero de Software, Motores de Bases de Datos}}
      {\tr{2022 -- present}{2022 -- presente}}
      \resumeItemListStart
        \resumeItem{\tr
          {\textbf{ColibriDB (columnar SQL engine, C++):} designed and built an embedded, vectorized engine on
           Apache Arrow that queries Parquet on local disk or remote storage over HTTP, from the SQL parser and
           optimizer to multi-threaded, push-based execution.}
          {\textbf{ColibriDB (motor SQL columnar, C++):} diseñé y construí un motor embedded y vectorizado sobre
           Apache Arrow que consulta archivos Parquet en disco local o en almacenamiento remoto vía HTTP, desde el
           parser SQL y el optimizador hasta la ejecución push-based multi-thread.}}
        \resumeItem{\tr
          {Compiled it to WebAssembly and shipped it as a TypeScript npm package that runs SQL in the browser,
           streaming results as Arrow record batches.}
          {Lo compilé a WebAssembly y lo publiqué como un paquete npm en TypeScript que ejecuta SQL en el
           navegador y entrega los resultados como record batches de Arrow.}}
        \resumeItem{\tr
          {\textbf{MillenniumDB (open-source multi-model graph database, C++):} built SPARQL 1.1 end to end,
           from the grammar to worst-case optimal join plans and operators for expressions, property paths and
           updates. The engine serves a public Wikidata endpoint and ran 3--10$\times$ faster than
           Virtuoso, Blazegraph, Jena and Neo4j on 1.26B triples.}
          {\textbf{MillenniumDB (base de datos de grafos multimodelo open source, C++):} implementé SPARQL 1.1
           completo, desde la gramática hasta planes con joins worst-case optimal y operadores de expresiones,
           property paths y updates. Sirve un endpoint público de Wikidata y fue 3--10$\times$ más rápido que
           Virtuoso, Blazegraph, Jena y Neo4j en 1.260M de triples.}}
        \resumeItem{\tr
          {Built vector search into the engine: HNSW and LSH indexes with online updates, SIMD distance kernels, and
           similarity-ordered iterators that plug into the join. Published the hybrid graph + vector algorithms at
           ISWC 2025: up to 50$\times$ faster than evaluate-then-rank on queries with many results.}
          {Incorporé búsqueda vectorial al motor: índices HNSW y LSH con actualizaciones en línea, kernels de
           distancia SIMD e iteradores ordenados por similitud que se integran al join. Publiqué los algoritmos
           híbridos grafo + vector en ISWC 2025: hasta 50$\times$ más rápidos que evaluar y luego ordenar en
           consultas con muchos resultados.}}
        \resumeItem{\tr
          {Built the full-text index, tensor storage for embeddings, an async Boost.Asio/Beast server with a
           chunked binary protocol over TCP and WebSocket, and the official Python (PyPI) and JavaScript (npm) drivers.}
          {Construí el full-text index, un almacén de tensores para embeddings, un servidor asíncrono
           Boost.Asio/Beast con protocolo binario vía TCP/WebSocket y los drivers oficiales de Python y JavaScript
           (PyPI, npm).}}
        \resumeItem{\tr
          {Mentored and supervised interns: onboarding, project scoping, and code reviews through to merge.}
          {Guié y supervisé practicantes: onboarding, alcance de sus proyectos y code reviews hasta el merge.}}
      \resumeItemListEnd

    \resumeSubheading
      {Pontificia Universidad Católica de Chile}
      {Santiago, Chile}
      {\tr{Instructor, Data Management for Public Policy}
          {Profesor, Gestión de Datos para Políticas Públicas}}
      {\tr{2024 -- present}{2024 -- presente}}
      \resumeItemListStart
        \resumeItem{\tr
          {Designed the course and teach it (relational modeling, SQL, R) to cohorts of 50+ students; 200+ to date.}
          {Diseñé el curso y lo dicto (modelo relacional, SQL, R) a cursos de más de 50 estudiantes; más de 200
           a la fecha.}}
      \resumeItemListEnd
    \resumeSubSubheading
      {\tr{Lecturer, Database Systems (IIC2413)}
          {Profesor, Bases de Datos (IIC2413)}}
      {2023}
      \resumeItemListStart
        \resumeItem{\tr
          {Taught relational algebra, SQL, indexing and NoSQL databases to 100+ undergraduates.}
          {Enseñé álgebra relacional, SQL, índices y bases de datos NoSQL a más de 100 estudiantes de pregrado.}}
      \resumeItemListEnd

    \resumeSubheading
      {coreDevX}{Santiago, Chile}
      {\tr{Software Engineer Intern}
          {Práctica en Ingeniería de Software}}
      {2021}
      \resumeItemListStart
        \resumeItem{\tr
          {Built a Node.js dependency-analysis tool for a large C\#, Java and PL/SQL codebase: it mapped cross-file
           dependencies into an ArangoDB graph, so teams could see which files they shared and coordinate changes.}
          {Construí en Node.js una herramienta de análisis de dependencias para un código extenso en C\#, Java y
           PL/SQL, que llevaba las dependencias entre archivos a un grafo en ArangoDB para coordinar cambios.}}
      \resumeItemListEnd

  \resumeSubHeadingListEnd

%----------OPEN SOURCE----------
\section{\tr{Open Source}{Código Abierto}}
    \resumeSubHeadingListStart
      \resumeProjectHeading
          {\textbf{\href{https://github.com/REmatchChile/REmatch}{REmatch}} $|$
           \emph{\tr{Regex engine that enumerates all matches}{motor de expresiones regulares que enumera todas las coincidencias}}}
          {2018 -- \tr{present}{presente}}
          \resumeItemListStart
            \resumeItem{\tr
              {Core C++ contributor; owned the build (CMake, Python wheels with pybind11, WebAssembly) and built
               REmatch.js and the \href{https://rematch.cl}{\underline{rematch.cl}} website.}
              {Contribuí al motor en C++; me encargué del build (CMake, wheels de Python con pybind11,
               WebAssembly) y desarrollé REmatch.js y el sitio \href{https://rematch.cl}{\underline{rematch.cl}}.}}
          \resumeItemListEnd
      \resumeProjectHeading
          {\textbf{\href{https://github.com/CORE-cer/CORE}{CORE+}} $|$
           \emph{\tr{Complex event recognition engine}{motor de reconocimiento de eventos complejos}}}
          {2025 -- 2026}
          \resumeItemListStart
            \resumeItem{\tr
              {Built the web client and the WebSocket layer that streams live query results.}
              {Construí el cliente web y la capa WebSocket que transmite en vivo los resultados de las consultas.}}
          \resumeItemListEnd
    \resumeSubHeadingListEnd

%----------PUBLICATIONS----------
\section{\tr{Publications}{Publicaciones}}
    \resumeCompactListStart
      \resumePublication
        {10.1007/978-3-032-09527-5_19}
        {Graph Querying or Similarity Search? Both!}{ISWC 2025}
      \resumePublication
        {10.1145/3722212.3725090}
        {CORE+: A Complex Event Recognition Engine in C++}{SIGMOD 2025}
      \resumePublication
        {10.1145/3626246.3654757}
        {MillenniumDB: A Multi-modal, Multi-model Graph Database}{SIGMOD 2024}
      \resumePublication
        {10.1145/3626246.3654746}
        {Demonstrating REmatch: A Novel RegEx Engine for Finding all Matches}{SIGMOD 2024}
      \resumePublication
        {10.1145/3661304.3661899}
        {TelarKG: a Knowledge Graph of Chile's Constitutional Process}{GRADES-NDA 2024}
    \resumeCompactListEnd

%----------EDUCATION----------
\section{\tr{Education}{Educación}}
  \resumeSubHeadingListStart
    \resumeSubheading
      {Pontificia Universidad Católica de Chile}{Santiago, Chile}
      {\tr{BSc in Computer Science, Minor in Data Science. Thesis: Compiling SPARQL in MillenniumDB}
          {Lic. en Ciencia de la Computación, Minor en Ciencia de Datos. Tesis: SPARQL en MillenniumDB}}
      {2017 -- 2022}
  \resumeSubHeadingListEnd

%----------SKILLS----------
\section{\tr{Technical Skills}{Habilidades Técnicas}}
 \resumeCompactListStart
    \small{\item{
     \textbf{\tr{Query engines}{Motores de consulta}}{: \tr
       {parsing (ANTLR4), binding, planning, query optimization, vectorized and parallel execution}
       {parsing (ANTLR4), binding, planificación, optimización, ejecución vectorizada y paralela}} \\
     \textbf{\tr{Storage \& data}{Almacenamiento}}{: \tr
       {Apache Arrow, Parquet, buffer manager, B+ trees, vector indexes, RDF, property graphs}
       {Apache Arrow, Parquet, buffer manager, B+ trees, índices vectoriales, RDF, property graphs}} \\
     \textbf{\tr{Programming languages}{Lenguajes de programación}}{: C++, Python, TypeScript, JavaScript (React, Node.js), R} \\
     \textbf{\tr{Query languages}{Lenguajes de consulta}}{: SQL, SPARQL, Cypher, GQL} \\
     \textbf{\tr{Tooling}{Herramientas}}{: CMake, Emscripten, WebAssembly, pybind11, Boost, GoogleTest,
       sanitizers, perf, Valgrind, gdb, Git, Docker, GitHub~Actions, Linux} \\
     \textbf{\tr{Spoken languages}{Idiomas}}{: \tr{Spanish (native), English (advanced)}{español (nativo), inglés (avanzado)}}
    }}
 \resumeCompactListEnd

\end{document}
